\documentclass[11pt,a4paper]{article}
\usepackage[utf8]{inputenc}
\usepackage[T1]{fontenc}
\usepackage[spanish,es-minimal,es-noshorthands]{babel}
\usepackage{lmodern}
\usepackage[sfdefault,lining]{FiraSans}
\usepackage[a4paper,top=2.3cm,bottom=2.2cm,left=1.8cm,right=1.8cm,headheight=1.3cm,headsep=0.4cm,footskip=1cm]{geometry}
\usepackage{amsmath,amssymb}
\usepackage{icomma}
\usepackage{graphicx}
\usepackage[table]{xcolor}
\usepackage{tikz}
\usepackage{fontawesome5}
\usepackage{ccicons}
\usepackage{fancyhdr}
\usepackage{lastpage}
\usepackage{tasks}
\usepackage{needspace}
\usepackage[most]{tcolorbox}
\usepackage{array}
\usepackage[hidelinks]{hyperref}
\AtBeginDocument{\hypersetup{pdftitle={Ficha de ejercicios de la unidad 10: Proporcionalidad. Teorema de Tales. 3º ESO},pdfsubject={Matemáticas 3º ESO},pdfkeywords={razón, proporción, segmentos proporcionales, teorema de Tales, posición de Tales, semejanza, razón de semejanza, criterios de semejanza, escala}}}

\definecolor{navy}{HTML}{1B2748}
\definecolor{accent}{HTML}{E8603A}
\definecolor{teal}{HTML}{1F9E90}
\definecolor{soft}{HTML}{EEF1F8}
\definecolor{muted}{HTML}{5D6685}

\setlength{\parindent}{0pt}
\setlength{\parskip}{3pt}
\settasks{label=\alph*), label-format=\bfseries\color{navy}, label-width=1.3em, item-indent=1.8em, column-sep=0.6em, after-item-skip=5pt, before-skip=3pt, after-skip=4pt}

% ---------- capçalera i peu ----------
\newcommand{\llicencia}{\href{https://creativecommons.org/licenses/by-nc-sa/4.0/deed.es}{\ccbyncsa}\ \ Licencia Creative Commons Atribución-NoComercial-CompartirIgual 4.0 Internacional (CC BY-NC-SA 4.0)}
\pagestyle{fancy}
\fancyhf{}
\fancyhead[L]{\raisebox{-0.35\height}{\includegraphics[height=0.95cm]{logo-gva.png}}}
\fancyhead[C]{\small\color{muted}Ficha de ejercicios · Unidad 10: Proporcionalidad. Teorema de Tales · 3\textsuperscript{o} ESO}
\fancyhead[R]{\raisebox{-0.35\height}{\includegraphics[height=0.95cm]{logo-vermellar.png}}}
\fancyfoot[L]{\scriptsize\color{muted}\llicencia}
\fancyfoot[R]{\small\color{muted}\thepage}
\renewcommand{\headrulewidth}{0.4pt}
\renewcommand{\headrule}{\color{muted!40}\hrule height \headrulewidth}
\fancypagestyle{primera}{\fancyhf{}\renewcommand{\headrulewidth}{0pt}
  \fancyfoot[L]{\scriptsize\color{muted}\llicencia}
  \fancyfoot[R]{\small\color{muted}\thepage}}

% ---------- nivells i exercicis ----------
\newcommand{\estrelles}[1]{{\color{accent}\ifcase#1\or\faStar\faStar[regular]\faStar[regular]\or\faStar\faStar\faStar[regular]\or\faStar\faStar\faStar\fi}}
\newcounter{exer}
\newcommand{\ex}[1]{\par\needspace{4\baselineskip}\medskip\stepcounter{exer}\noindent{\color{navy}\bfseries\theexer.}\ \ {\small\estrelles{#1}}\ \ }

\newcommand{\seccio}[2]{\par\needspace{10\baselineskip}\bigskip
  \begin{tcolorbox}[colback=navy,colframe=navy,coltext=white,arc=2mm,boxrule=0pt,left=3mm,right=3mm,top=1.2mm,bottom=1.2mm]
  {\large\bfseries #1\quad #2}
  \end{tcolorbox}\vspace{-2pt}}

\newcolumntype{C}[1]{>{\centering\arraybackslash}p{#1}}
\usepackage{pifont}

% quadre de recordatori (només a la fitxa d'exercicis)
\newcommand{\recorda}[1]{\par\smallskip\begin{tcolorbox}[colback=accent!7,colframe=accent!55,arc=2mm,boxrule=0.5pt,left=2.5mm,right=2.5mm,top=1mm,bottom=1mm]\small{\color{accent}\faLightbulb[regular]}\ \textbf{Recuerda:}\ #1\end{tcolorbox}}
\newcommand{\res}[1]{#1}
\newcommand{\rest}[1]{#1}

\def\fmneed{14}
% =====================================================================
% Capçalera FentMates de les fitxes: un color per a cada curs.
% Els colors són els de la pàgina inicial (LEVELS a assets/js/materials.js);
% el logo ve de _eines/logo/pdf/ (python3 _eines/logo/genera.py).
%
% Ús, al final del preàmbul de la unitat:
%   % =====================================================================
% Capçalera FentMates de les fitxes: un color per a cada curs.
% Els colors són els de la pàgina inicial (LEVELS a assets/js/materials.js);
% el logo ve de _eines/logo/pdf/ (python3 _eines/logo/genera.py).
%
% Ús, al final del preàmbul de la unitat:
%   % =====================================================================
% Capçalera FentMates de les fitxes: un color per a cada curs.
% Els colors són els de la pàgina inicial (LEVELS a assets/js/materials.js);
% el logo ve de _eines/logo/pdf/ (python3 _eines/logo/genera.py).
%
% Ús, al final del preàmbul de la unitat:
%   \input{../comu/fentmates.tex}
%   \fmcurs{1eso}
% i, al començament del document:
%   \capcalera{Matemàtiques}{1r ESO}{Fitxa d'exercicis · Unitat 2}{Divisibilitat}
% Opcional abans de \input: \def\fmneed{14} (espai mínim abans de cada secció).
% =====================================================================
\makeatletter
\def\fm@logodir{../../logo/pdf}
\providecommand{\fmneed}{10}
\def\fm@cursnom{}

% clau del curs, color i color del text damunt del color
\newcommand{\fm@curs}[3]{\@namedef{fm@color@#1}{#2}\@namedef{fm@text@#1}{#3}}
\fm@curs{1eso}{E6B84A}{navy}
\fm@curs{2eso}{3B82F6}{white}
\fm@curs{3eso}{FF7A50}{white}
\fm@curs{4eso-a}{A78BFA}{white}
\fm@curs{4eso-b}{6D28D9}{white}
\fm@curs{1batx-mi}{9B1C1C}{white}
\fm@curs{1batx-ccss}{EF6B6B}{white}
\fm@curs{2batx-mii}{1E6B47}{white}
\fm@curs{2batx-ccss}{2BB3A3}{white}

\newcommand{\fmcurs}[1]{%
  \@ifundefined{fm@color@#1}{\PackageError{fentmates}{Curs desconegut: #1}{}}{}%
  \edef\fm@tmp{\noexpand\definecolor{curs}{HTML}{\@nameuse{fm@color@#1}}%
               \noexpand\colorlet{curstext}{\@nameuse{fm@text@#1}}}\fm@tmp
  \colorlet{cursfosc}{curs!72!black}%
  \def\fm@logo{\fm@logodir/vertical-#1.pdf}%
  \def\fm@marca{\fm@logodir/marca-#1.pdf}}

% ---------- capçalera de la primera pàgina ----------
% #1 matèria (petit), #2 curs, #3 tipus de document i unitat, #4 títol
\newcommand{\capcalera}[4]{%
  \gdef\fm@cursnom{#2}%
  \par\noindent
  \begin{tcolorbox}[colback=curs,colframe=curs,boxrule=0pt,arc=2.2mm,width=3.5cm,height=2.15cm,
      valign=center,left=3mm,right=2mm,top=1mm,bottom=1mm,nobeforeafter,box align=center]
    \color{curstext}{\footnotesize #1}\par\vspace{1pt}{\fontsize{19}{21}\selectfont\bfseries #2\par}%
  \end{tcolorbox}\hspace{5mm}%
  \begin{minipage}[c]{\dimexpr\linewidth-3.5cm-5mm-2.9cm\relax}
    {\color{cursfosc}\small\bfseries\MakeUppercase{#3}}\par\vspace{2pt}
    {\color{navy}\fontsize{20}{23}\selectfont\bfseries #4\par}
  \end{minipage}\hfill
  \begin{minipage}[c]{2.6cm}\raggedleft\includegraphics[width=2.6cm]{\fm@logo}\end{minipage}\par
  \vspace{2pt}}

% ---------- barres de secció del color del curs ----------
\renewcommand{\seccio}[2]{\par\needspace{\fmneed\baselineskip}\bigskip
  \begin{tcolorbox}[colback=curs,colframe=curs,coltext=curstext,arc=2mm,boxrule=0pt,left=3mm,right=3mm,top=1.2mm,bottom=1.2mm]
  {\large\bfseries #1\quad #2}
  \end{tcolorbox}\vspace{-2pt}}

% ---------- estrelles de dificultat del color del curs (a Batxillerat no se'n posen) ----------
\@ifundefined{estrelles}{}{%
  \renewcommand{\estrelles}[1]{{\color{curs}\ifcase#1\or\faStar\faStar[regular]\faStar[regular]\or\faStar\faStar\faStar[regular]\or\faStar\faStar\faStar\fi}}}

% ---------- capçalera i peu de les pàgines ----------
% primera pàgina: logos institucionals al peu
\fancypagestyle{primera}{\fancyhf{}\renewcommand{\headrulewidth}{0pt}
  \fancyfoot[L]{\raisebox{-0.32\height}{\includegraphics[height=0.78cm]{logo-gva.png}}\hspace{5mm}%
    \raisebox{-0.32\height}{\includegraphics[height=0.78cm]{logo-vermellar.png}}}
  \fancyfoot[R]{\scriptsize\color{muted}\href{https://creativecommons.org/licenses/by-nc-sa/4.0/deed.ca}{\ccbyncsa}\ \ CC BY-NC-SA 4.0 · fentmates.org\qquad\small\thepage}}
% la resta: el símbol a l'esquerra i el curs a la dreta, amb el títol al centre com abans
\fancyhead[L]{\raisebox{-0.3\height}{\includegraphics[height=0.62cm]{\fm@marca}}}
\fancyhead[R]{\small\bfseries\color{cursfosc}\fm@cursnom}
\renewcommand{\headrule}{\color{curs!45}\hrule height \headrulewidth}
\makeatother

%   \fmcurs{1eso}
% i, al començament del document:
%   \capcalera{Matemàtiques}{1r ESO}{Fitxa d'exercicis · Unitat 2}{Divisibilitat}
% Opcional abans de \input: \def\fmneed{14} (espai mínim abans de cada secció).
% =====================================================================
\makeatletter
\def\fm@logodir{../../logo/pdf}
\providecommand{\fmneed}{10}
\def\fm@cursnom{}

% clau del curs, color i color del text damunt del color
\newcommand{\fm@curs}[3]{\@namedef{fm@color@#1}{#2}\@namedef{fm@text@#1}{#3}}
\fm@curs{1eso}{E6B84A}{navy}
\fm@curs{2eso}{3B82F6}{white}
\fm@curs{3eso}{FF7A50}{white}
\fm@curs{4eso-a}{A78BFA}{white}
\fm@curs{4eso-b}{6D28D9}{white}
\fm@curs{1batx-mi}{9B1C1C}{white}
\fm@curs{1batx-ccss}{EF6B6B}{white}
\fm@curs{2batx-mii}{1E6B47}{white}
\fm@curs{2batx-ccss}{2BB3A3}{white}

\newcommand{\fmcurs}[1]{%
  \@ifundefined{fm@color@#1}{\PackageError{fentmates}{Curs desconegut: #1}{}}{}%
  \edef\fm@tmp{\noexpand\definecolor{curs}{HTML}{\@nameuse{fm@color@#1}}%
               \noexpand\colorlet{curstext}{\@nameuse{fm@text@#1}}}\fm@tmp
  \colorlet{cursfosc}{curs!72!black}%
  \def\fm@logo{\fm@logodir/vertical-#1.pdf}%
  \def\fm@marca{\fm@logodir/marca-#1.pdf}}

% ---------- capçalera de la primera pàgina ----------
% #1 matèria (petit), #2 curs, #3 tipus de document i unitat, #4 títol
\newcommand{\capcalera}[4]{%
  \gdef\fm@cursnom{#2}%
  \par\noindent
  \begin{tcolorbox}[colback=curs,colframe=curs,boxrule=0pt,arc=2.2mm,width=3.5cm,height=2.15cm,
      valign=center,left=3mm,right=2mm,top=1mm,bottom=1mm,nobeforeafter,box align=center]
    \color{curstext}{\footnotesize #1}\par\vspace{1pt}{\fontsize{19}{21}\selectfont\bfseries #2\par}%
  \end{tcolorbox}\hspace{5mm}%
  \begin{minipage}[c]{\dimexpr\linewidth-3.5cm-5mm-2.9cm\relax}
    {\color{cursfosc}\small\bfseries\MakeUppercase{#3}}\par\vspace{2pt}
    {\color{navy}\fontsize{20}{23}\selectfont\bfseries #4\par}
  \end{minipage}\hfill
  \begin{minipage}[c]{2.6cm}\raggedleft\includegraphics[width=2.6cm]{\fm@logo}\end{minipage}\par
  \vspace{2pt}}

% ---------- barres de secció del color del curs ----------
\renewcommand{\seccio}[2]{\par\needspace{\fmneed\baselineskip}\bigskip
  \begin{tcolorbox}[colback=curs,colframe=curs,coltext=curstext,arc=2mm,boxrule=0pt,left=3mm,right=3mm,top=1.2mm,bottom=1.2mm]
  {\large\bfseries #1\quad #2}
  \end{tcolorbox}\vspace{-2pt}}

% ---------- estrelles de dificultat del color del curs (a Batxillerat no se'n posen) ----------
\@ifundefined{estrelles}{}{%
  \renewcommand{\estrelles}[1]{{\color{curs}\ifcase#1\or\faStar\faStar[regular]\faStar[regular]\or\faStar\faStar\faStar[regular]\or\faStar\faStar\faStar\fi}}}

% ---------- capçalera i peu de les pàgines ----------
% primera pàgina: logos institucionals al peu
\fancypagestyle{primera}{\fancyhf{}\renewcommand{\headrulewidth}{0pt}
  \fancyfoot[L]{\raisebox{-0.32\height}{\includegraphics[height=0.78cm]{logo-gva.png}}\hspace{5mm}%
    \raisebox{-0.32\height}{\includegraphics[height=0.78cm]{logo-vermellar.png}}}
  \fancyfoot[R]{\scriptsize\color{muted}\href{https://creativecommons.org/licenses/by-nc-sa/4.0/deed.ca}{\ccbyncsa}\ \ CC BY-NC-SA 4.0 · fentmates.org\qquad\small\thepage}}
% la resta: el símbol a l'esquerra i el curs a la dreta, amb el títol al centre com abans
\fancyhead[L]{\raisebox{-0.3\height}{\includegraphics[height=0.62cm]{\fm@marca}}}
\fancyhead[R]{\small\bfseries\color{cursfosc}\fm@cursnom}
\renewcommand{\headrule}{\color{curs!45}\hrule height \headrulewidth}
\makeatother

%   \fmcurs{1eso}
% i, al començament del document:
%   \capcalera{Matemàtiques}{1r ESO}{Fitxa d'exercicis · Unitat 2}{Divisibilitat}
% Opcional abans de \input: \def\fmneed{14} (espai mínim abans de cada secció).
% =====================================================================
\makeatletter
\def\fm@logodir{../../logo/pdf}
\providecommand{\fmneed}{10}
\def\fm@cursnom{}

% clau del curs, color i color del text damunt del color
\newcommand{\fm@curs}[3]{\@namedef{fm@color@#1}{#2}\@namedef{fm@text@#1}{#3}}
\fm@curs{1eso}{E6B84A}{navy}
\fm@curs{2eso}{3B82F6}{white}
\fm@curs{3eso}{FF7A50}{white}
\fm@curs{4eso-a}{A78BFA}{white}
\fm@curs{4eso-b}{6D28D9}{white}
\fm@curs{1batx-mi}{9B1C1C}{white}
\fm@curs{1batx-ccss}{EF6B6B}{white}
\fm@curs{2batx-mii}{1E6B47}{white}
\fm@curs{2batx-ccss}{2BB3A3}{white}

\newcommand{\fmcurs}[1]{%
  \@ifundefined{fm@color@#1}{\PackageError{fentmates}{Curs desconegut: #1}{}}{}%
  \edef\fm@tmp{\noexpand\definecolor{curs}{HTML}{\@nameuse{fm@color@#1}}%
               \noexpand\colorlet{curstext}{\@nameuse{fm@text@#1}}}\fm@tmp
  \colorlet{cursfosc}{curs!72!black}%
  \def\fm@logo{\fm@logodir/vertical-#1.pdf}%
  \def\fm@marca{\fm@logodir/marca-#1.pdf}}

% ---------- capçalera de la primera pàgina ----------
% #1 matèria (petit), #2 curs, #3 tipus de document i unitat, #4 títol
\newcommand{\capcalera}[4]{%
  \gdef\fm@cursnom{#2}%
  \par\noindent
  \begin{tcolorbox}[colback=curs,colframe=curs,boxrule=0pt,arc=2.2mm,width=3.5cm,height=2.15cm,
      valign=center,left=3mm,right=2mm,top=1mm,bottom=1mm,nobeforeafter,box align=center]
    \color{curstext}{\footnotesize #1}\par\vspace{1pt}{\fontsize{19}{21}\selectfont\bfseries #2\par}%
  \end{tcolorbox}\hspace{5mm}%
  \begin{minipage}[c]{\dimexpr\linewidth-3.5cm-5mm-2.9cm\relax}
    {\color{cursfosc}\small\bfseries\MakeUppercase{#3}}\par\vspace{2pt}
    {\color{navy}\fontsize{20}{23}\selectfont\bfseries #4\par}
  \end{minipage}\hfill
  \begin{minipage}[c]{2.6cm}\raggedleft\includegraphics[width=2.6cm]{\fm@logo}\end{minipage}\par
  \vspace{2pt}}

% ---------- barres de secció del color del curs ----------
\renewcommand{\seccio}[2]{\par\needspace{\fmneed\baselineskip}\bigskip
  \begin{tcolorbox}[colback=curs,colframe=curs,coltext=curstext,arc=2mm,boxrule=0pt,left=3mm,right=3mm,top=1.2mm,bottom=1.2mm]
  {\large\bfseries #1\quad #2}
  \end{tcolorbox}\vspace{-2pt}}

% ---------- estrelles de dificultat del color del curs (a Batxillerat no se'n posen) ----------
\@ifundefined{estrelles}{}{%
  \renewcommand{\estrelles}[1]{{\color{curs}\ifcase#1\or\faStar\faStar[regular]\faStar[regular]\or\faStar\faStar\faStar[regular]\or\faStar\faStar\faStar\fi}}}

% ---------- capçalera i peu de les pàgines ----------
% primera pàgina: logos institucionals al peu
\fancypagestyle{primera}{\fancyhf{}\renewcommand{\headrulewidth}{0pt}
  \fancyfoot[L]{\raisebox{-0.32\height}{\includegraphics[height=0.78cm]{logo-gva.png}}\hspace{5mm}%
    \raisebox{-0.32\height}{\includegraphics[height=0.78cm]{logo-vermellar.png}}}
  \fancyfoot[R]{\scriptsize\color{muted}\href{https://creativecommons.org/licenses/by-nc-sa/4.0/deed.ca}{\ccbyncsa}\ \ CC BY-NC-SA 4.0 · fentmates.org\qquad\small\thepage}}
% la resta: el símbol a l'esquerra i el curs a la dreta, amb el títol al centre com abans
\fancyhead[L]{\raisebox{-0.3\height}{\includegraphics[height=0.62cm]{\fm@marca}}}
\fancyhead[R]{\small\bfseries\color{cursfosc}\fm@cursnom}
\renewcommand{\headrule}{\color{curs!45}\hrule height \headrulewidth}
\makeatother

\fmcurs{3eso}

\begin{document}
\thispagestyle{primera}

\capcalera{Matemáticas}{3\textsuperscript{o} ESO}{Ficha de ejercicios · Unidad 10}{Proporcionalidad. Teorema de Tales}
\vspace{10pt}
Nombre y apellidos: \rule{6.6cm}{0.4pt}\hfill Grupo: \rule{1.6cm}{0.4pt}\hfill Fecha: \rule{2.2cm}{0.4pt}
\vspace{8pt}

\begin{tcolorbox}[colback=soft,colframe=soft,arc=2mm,left=3mm,right=3mm,top=1.5mm,bottom=1.5mm]
\small Ejercicios graduados por niveles de dificultad:\quad \estrelles{1}\ básico \qquad \estrelles{2}\ medio \qquad \estrelles{3}\ avanzado.\\[2pt] Resuélvelos en el cuaderno mostrando todos los pasos y haz un dibujo de cada situación. Escribe la proporción antes de calcular y, si un resultado no es exacto, redondéalo a las centésimas.
\end{tcolorbox}

% =====================================================
\seccio{1}{Razones y proporciones entre segmentos}

\ex{1} Calcula la razón entre los segmentos (si no es exacta, redondea a las centésimas):\begin{tasks}(4)
  \task 12 cm y 8 cm
  \task 5 cm y 20 cm
  \task 0,6 m y 45 cm
  \task 2 dm y 25 cm
\end{tasks}\recorda{La \textbf{razón} de dos segmentos es el cociente de sus longitudes, medidas con la \textbf{misma unidad}. Cuatro segmentos son \textbf{proporcionales} si $\frac ab=\frac cd$, es decir, si $a\cdot d=b\cdot c$ (producto de los extremos = producto de los medios).}

\ex{1} Calcula $x$ en cada proporción:\begin{tasks}(3)
  \task $\dfrac35=\dfrac{x}{20}$
  \task $\dfrac8x=\dfrac27$
  \task $\dfrac x4=\dfrac{15}{6}$
  \task $\dfrac{2{,}5}{x}=\dfrac58$
  \task $\dfrac69=\dfrac{10}{x}$
  \task $\dfrac{x}{12}=\dfrac{0{,}75}{3}$
\end{tasks}

\ex{1} ¿Son proporcionales, en este orden, los siguientes segmentos?\begin{tasks}(2)
  \task 2, 3, 4 y 6 cm
  \task 3, 5, 6 y 9 cm
  \task 1,5, 2, 4,5 y 6 cm
  \task 4, 7, 8 y 15 cm
\end{tasks}

\ex{2} Calcula el cuarto proporcional de los segmentos: \ a) 4, 6 y 10 cm; \ b) 2,5, 3 y 7 cm.

\ex{2} Dos segmentos están en razón $\frac35$. Si el pequeño mide 12 cm, ¿cuánto mide el grande?

\ex{2} Dos segmentos suman 56 cm y están en razón $\frac34$. ¿Cuánto mide cada uno?

\ex{3} La \textbf{media proporcional} de $a$ y $b$ es el segmento $x$ tal que $\frac ax=\frac xb$. Calcula la media proporcional de: \ a) 4 y 25; \ b) 2 y 8; \ c) 3 y 6 (redondea a las centésimas).


% =====================================================
\seccio{2}{El teorema de Tales}

\ex{1} Las rectas \textcolor{navy}{azul oscuro} son paralelas. Calcula $x$.\par\smallskip\begin{center}\begin{tikzpicture}[x=0.55cm,y=0.55cm,baseline=(current bounding box.center),line join=round]
\draw[draw=navy,thick] (-1.200,0.000) -- (5.503,0.000);
\draw[draw=navy,thick] (-1.200,2.000) -- (5.503,2.000);
\draw[draw=navy,thick] (-1.200,5.000) -- (5.503,5.000);
\draw[thick] (-0.490,-0.700) -- (3.991,5.700);
\draw[thick] (3.502,-0.700) -- (4.401,5.700);
\draw[draw=accent,line width=1.6pt] (3.881,2.000) -- (4.303,5.000);
\fill (0.000,0.000) circle (1.5pt);
\fill (1.400,2.000) circle (1.5pt);
\fill (3.501,5.000) circle (1.5pt);
\fill (3.600,0.000) circle (1.5pt);
\fill (3.881,2.000) circle (1.5pt);
\fill (4.303,5.000) circle (1.5pt);
\node[font=\small,inner sep=2pt,above left] at (0.000,0.000) {$A$};
\node[font=\small,inner sep=2pt,above left] at (1.400,2.000) {$B$};
\node[font=\small,inner sep=2pt,above left] at (3.501,5.000) {$C$};
\node[font=\small,inner sep=2pt,above right] at (3.600,0.000) {$A'$};
\node[font=\small,inner sep=2pt,above right] at (3.881,2.000) {$B'$};
\node[font=\small,inner sep=2pt,above right] at (4.303,5.000) {$C'$};
\node[font=\small,inner sep=2pt,left] at (0.700,1.000) {$4$};
\node[font=\small,inner sep=2pt,left] at (2.451,3.500) {$6$};
\node[font=\small,inner sep=2pt,right] at (3.741,1.000) {$6$};
\node[font=\small,inner sep=2pt,right,text=accent] at (4.092,3.500) {$x$};
\end{tikzpicture}\end{center}\recorda{\textbf{Teorema de Tales:} si unas rectas paralelas cortan dos secantes, los segmentos que determinan sobre una son proporcionales a los de la otra: $\frac{\overline{AB}}{\overline{BC}}=\frac{\overline{A'B'}}{\overline{B'C'}}$, o bien $\frac{\overline{AB}}{\overline{A'B'}}=\frac{\overline{BC}}{\overline{B'C'}}=\frac{\overline{AC}}{\overline{A'C'}}$.}

\ex{1} Las rectas $r$, $s$ y $t$ son paralelas y cortan las secantes en $A,\,B,\,C$ y $A',\,B',\,C'$. Calcula el segmento que falta:\begin{tasks}(1)
  \task $\overline{AB}=3$, $\overline{BC}=5$, $\overline{A'B'}=4{,}8$. Calcula $\overline{B'C'}$.
  \task $\overline{AB}=6$, $\overline{A'B'}=4$, $\overline{B'C'}=5$. Calcula $\overline{BC}$.
  \task $\overline{AB}=4$, $\overline{AC}=10$, $\overline{A'C'}=15$. Calcula $\overline{A'B'}$.
\end{tasks}

\ex{2} Las rectas \textcolor{navy}{azul oscuro} son paralelas. Calcula $x$.\par\smallskip\begin{center}\begin{tikzpicture}[x=0.55cm,y=0.55cm,baseline=(current bounding box.center),line join=round]
\draw[draw=navy,thick] (-1.200,0.000) -- (5.475,0.000);
\draw[draw=navy,thick] (-1.200,1.600) -- (5.475,1.600);
\draw[draw=navy,thick] (-1.200,4.800) -- (5.475,4.800);
\draw[thick] (-0.490,-0.700) -- (3.851,5.500);
\draw[thick] (3.502,-0.700) -- (4.373,5.500);
\draw[draw=accent,line width=1.6pt] (0.000,0.000) -- (1.120,1.600);
\fill (0.000,0.000) circle (1.5pt);
\fill (1.120,1.600) circle (1.5pt);
\fill (3.361,4.800) circle (1.5pt);
\fill (3.600,0.000) circle (1.5pt);
\fill (3.825,1.600) circle (1.5pt);
\fill (4.275,4.800) circle (1.5pt);
\node[font=\small,inner sep=2pt,above left] at (0.000,0.000) {$A$};
\node[font=\small,inner sep=2pt,above left] at (1.120,1.600) {$B$};
\node[font=\small,inner sep=2pt,above left] at (3.361,4.800) {$C$};
\node[font=\small,inner sep=2pt,above right] at (3.600,0.000) {$A'$};
\node[font=\small,inner sep=2pt,above right] at (3.825,1.600) {$B'$};
\node[font=\small,inner sep=2pt,above right] at (4.275,4.800) {$C'$};
\node[font=\small,inner sep=2pt,left,text=accent] at (0.560,0.800) {$x$};
\node[font=\small,inner sep=2pt,left] at (2.241,3.200) {$6$};
\node[font=\small,inner sep=2pt,right] at (3.712,0.800) {$4$};
\node[font=\small,inner sep=2pt,right] at (4.050,3.200) {$8$};
\end{tikzpicture}\end{center}

\ex{2} Las rectas que cortan dos secantes determinan en ellas estos segmentos. ¿Son paralelas?\begin{tasks}(3)
  \task 2 y 3; \ 4 y 6
  \task 3 y 4; \ 4 y 5
  \task 1,5 y 2,5; \ 2,4 y 4
\end{tasks}

\ex{2} Tres parcelas están entre dos calles y sus lindes son paralelas. En la calle del Mar miden 15 m, 20 m y 25 m de fachada. En la calle Ancha, las tres juntas miden 72 m. ¿Cuánto mide la fachada de cada parcela en la calle Ancha?

\ex{3} Unas rectas paralelas determinan $\overline{AB}=x$ y $\overline{BC}=x+4$ sobre una secante, y $\overline{A'B'}=3$ y $\overline{B'C'}=5$ sobre la otra. Calcula $x$.

\ex{3} Unas rectas paralelas determinan sobre una secante $\overline{AB}=5$ cm y $\overline{BC}=8$ cm. Sobre otra secante, $\overline{A'C'}=19{,}5$ cm. Calcula $\overline{A'B'}$ y $\overline{B'C'}$.


% =====================================================
\seccio{3}{Dividir un segmento en partes iguales o proporcionales}

\ex{1} Dibuja un segmento de 12 cm y divídelo en 5 partes iguales con regla y compás. ¿Cuánto mide cada parte?\recorda{Desde un extremo, traza una semirrecta auxiliar y marca en ella con el compás tantos segmentos iguales como partes quieras. Une la última marca con el otro extremo y traza paralelas por las otras marcas. Para partes proporcionales a $2,\,3,\,4$, marca $2+3+4=9$ segmentos y traza las paralelas por las marcas 2 y 5.}

\ex{1} Calcula las partes en que queda dividido:\begin{tasks}(1)
  \task un segmento de 30 cm, en partes proporcionales a 2 y 3
  \task un segmento de 24 cm, en partes proporcionales a 1, 2 y 3
  \task un segmento de 45 cm, en partes proporcionales a 2, 3 y 4
\end{tasks}

\ex{2} Dibuja un segmento de 10 cm y divídelo, con regla y compás, en dos partes proporcionales a 2 y 3. Explica los pasos y comprueba el resultado calculando.

\ex{2} Un punto $P$ divide un segmento $AB$ de 35 cm de manera que $\frac{\overline{AP}}{\overline{PB}}=\frac25$. ¿Cuánto miden $\overline{AP}$ y $\overline{PB}$?

\ex{2} Cortamos una cuerda de 3,6 m en tres trozos proporcionales a 1, 2 y 5. ¿Cuánto mide cada trozo?

\ex{3} Un punto $P$ divide un segmento $AB$ en dos partes que están en razón $\frac37$, y la parte larga mide 12 cm más que la corta. ¿Cuánto miden las dos partes y el segmento entero?


% =====================================================
\seccio{4}{Triángulos en posición de Tales}

\ex{1} En la figura, $B'C'\parallel BC$. Calcula $x$.\par\smallskip\begin{center}\begin{tikzpicture}[x=0.5cm,y=0.5cm,baseline=(current bounding box.center),line join=round]
\draw[fill=teal!12,draw=teal,thick] (0.000,0.000) -- (-2.817,-5.298) -- (2.644,-5.938) -- cycle;
\draw[fill=accent!22,draw=accent,thick] (0.000,0.000) -- (-1.127,-2.119) -- (1.058,-2.375) -- cycle;
\draw[draw=accent,line width=1.6pt] (-1.127,-2.119) -- (1.058,-2.375);
\draw[draw=teal,line width=1.6pt] (-2.817,-5.298) -- (2.644,-5.938);
\fill (0.000,0.000) circle (1.5pt);
\fill (-2.817,-5.298) circle (1.5pt);
\fill (2.644,-5.938) circle (1.5pt);
\fill (-1.127,-2.119) circle (1.5pt);
\fill (1.058,-2.375) circle (1.5pt);
\node[font=\small,inner sep=2pt,above] at (0.000,0.000) {$A$};
\node[font=\small,inner sep=2pt,left] at (-2.817,-5.298) {$B$};
\node[font=\small,inner sep=2pt,right] at (2.644,-5.938) {$C$};
\node[font=\small,inner sep=2pt,left] at (-1.127,-2.119) {$B'$};
\node[font=\small,inner sep=2pt,right] at (1.058,-2.375) {$C'$};
\node[font=\small,inner sep=2pt,left] at (-0.563,-1.060) {$4$};
\node[font=\small,inner sep=2pt,left] at (-1.972,-3.708) {$6$};
\node[font=\small,inner sep=2pt,below] at (-0.087,-5.618) {$15$};
\node[font=\small,inner sep=2pt,above,text=accent] at (-0.035,-2.247) {$x$};
\end{tikzpicture}\end{center}\recorda{Si $B'C'\parallel BC$, los triángulos $AB'C'$ y $ABC$ están en \textbf{posición de Tales}: $\frac{\overline{AB'}}{\overline{AB}}=\frac{\overline{AC'}}{\overline{AC}}=\frac{\overline{B'C'}}{\overline{BC}}$. Para el lado paralelo, compara siempre \textbf{lados enteros} desde el vértice $A$.}

\ex{1} En un triángulo $ABC$, $B'$ es un punto de $AB$ y $C'$ un punto de $AC$, con $B'C'\parallel BC$. Calcula:\begin{tasks}(1)
  \task $\overline{B'C'}$, si $\overline{AB'}=3$, $\overline{B'B}=6$ y $\overline{BC}=12$.
  \task $\overline{BC}$, si $\overline{AB'}=5$, $\overline{AB}=8$ y $\overline{B'C'}=7{,}5$.
  \task $\overline{BC}$, si $\overline{AC'}=4$, $\overline{C'C}=2$ y $\overline{B'C'}=6$.
\end{tasks}

\ex{2} En la figura, $B'C'\parallel BC$. Calcula $x$ y $y$.\par\smallskip\begin{center}\begin{tikzpicture}[x=0.5cm,y=0.5cm,baseline=(current bounding box.center),line join=round]
\draw[fill=teal!12,draw=teal,thick] (0.000,0.000) -- (-2.817,-5.298) -- (2.644,-5.938) -- cycle;
\draw[fill=accent!22,draw=accent,thick] (0.000,0.000) -- (-1.690,-3.179) -- (1.586,-3.563) -- cycle;
\draw[draw=accent,line width=1.6pt] (-1.690,-3.179) -- (1.586,-3.563);
\draw[draw=teal,line width=1.6pt] (-2.817,-5.298) -- (2.644,-5.938);
\fill (0.000,0.000) circle (1.5pt);
\fill (-2.817,-5.298) circle (1.5pt);
\fill (2.644,-5.938) circle (1.5pt);
\fill (-1.690,-3.179) circle (1.5pt);
\fill (1.586,-3.563) circle (1.5pt);
\node[font=\small,inner sep=2pt,above] at (0.000,0.000) {$A$};
\node[font=\small,inner sep=2pt,left] at (-2.817,-5.298) {$B$};
\node[font=\small,inner sep=2pt,right] at (2.644,-5.938) {$C$};
\node[font=\small,inner sep=2pt,left] at (-1.690,-3.179) {$B'$};
\node[font=\small,inner sep=2pt,right] at (1.586,-3.563) {$C'$};
\node[font=\small,inner sep=2pt,left] at (-0.845,-1.589) {$6$};
\node[font=\small,inner sep=2pt,left] at (-2.253,-4.238) {$4$};
\node[font=\small,inner sep=2pt,right] at (0.793,-1.781) {$9$};
\node[font=\small,inner sep=2pt,right] at (2.115,-4.750) {$y$};
\node[font=\small,inner sep=2pt,above] at (-0.052,-3.371) {$7{,}2$};
\node[font=\small,inner sep=2pt,below] at (-0.087,-5.618) {$x$};
\end{tikzpicture}\end{center}

\ex{2} Un triángulo tiene 18 cm de base y 12 cm de altura. Trazamos una paralela a la base a 4 cm del vértice opuesto. ¿Cuánto mide el segmento de la paralela que queda dentro del triángulo?

\ex{2} En la figura, $AB\parallel DE$. Calcula $x=\overline{DE}$ y $y=\overline{CD}$.\par\smallskip\begin{center}\begin{tikzpicture}[x=0.42cm,y=0.42cm,baseline=(current bounding box.center),line join=round]
\draw[fill=teal!12,draw=teal,thick] (0.000,0.000) -- (6.000,0.000) -- (2.400,3.200) -- cycle;
\draw[fill=accent!22,draw=accent,thick] (2.400,3.200) -- (6.000,8.000) -- (-3.000,8.000) -- cycle;
\node[font=\small,below left] at (0.000,0.000) {$A$};
\node[font=\small,below right] at (6.000,0.000) {$B$};
\node[font=\small,left=4pt] at (2.400,3.200) {$C$};
\node[font=\small,above left] at (-3.000,8.000) {$D$};
\node[font=\small,above right] at (6.000,8.000) {$E$};
\node[font=\small,below] at (3.000,0.000) {$8$};
\node[font=\small,left] at (1.200,1.600) {$6$};
\node[font=\small,right] at (4.200,1.600) {$5$};
\node[font=\small,above] at (1.500,8.000) {$x$};
\node[font=\small,right] at (4.200,5.600) {$9$};
\node[font=\small,left] at (-0.300,5.600) {$y$};
\end{tikzpicture}\end{center}

\ex{3} Para medir la anchura de un río, Álex clava una estaca $B$ en la orilla, delante de un árbol $A$ de la otra orilla. Camina 15 m perpendicularmente a $AB$ hasta $C$, clava otra, continúa 5 m hasta $D$ y gira 90° alejándose del río hasta que ve $C$ y $A$ alineados, en el punto $E$, que está a 4 m de $D$. ¿Qué anchura tiene el río?

\ex{3} En un triángulo $ABC$, $B'C'\parallel BC$, $\overline{B'B}=6$, $\overline{B'C'}=4$ y $\overline{BC}=10$. Calcula $x=\overline{AB'}$.


% =====================================================
\seccio{5}{Figuras semejantes. Razón de semejanza}

\ex{1} ¿Son semejantes estas parejas de rectángulos? Si lo son, calcula la razón de semejanza.\begin{tasks}(2)
  \task $4\times6$ \ y \ $6\times9$
  \task $2\times5$ \ y \ $4\times8$
  \task $3\times4$ \ y \ $6\times8$
  \task un cuadrado de 3 cm y uno de 7 cm
\end{tasks}\recorda{Dos figuras son \textbf{semejantes} si tienen los ángulos iguales y los lados homólogos proporcionales. La \textbf{razón de semejanza} es $k=\frac{a'}{a}$. Los perímetros están en razón $k$ y las áreas en razón $k^2$.}

\ex{1} Un triángulo tiene lados de 4, 6 y 7 cm. Calcula los lados de un triángulo semejante con razón: \ a) $k=2{,}5$; \ b) $k=0{,}4$.

\ex{2} Una foto de 9 cm × 13 cm se amplía hasta que el lado pequeño mide 27 cm. ¿Cuánto mide el lado grande? ¿Cuántas veces mayor es el área?

\ex{2} Dos pentágonos son semejantes con razón $k=1{,}5$. El pequeño tiene 20 cm de perímetro y 27,5 cm² de área. Calcula el perímetro y el área del grande.

\ex{2} Dos figuras semejantes tienen áreas de 18 cm² y 72 cm². ¿Cuál es la razón de semejanza? Si el perímetro de la pequeña es 16 cm, ¿cuál es el de la grande?

\ex{3} Tenemos un rectángulo de 8 cm × 12 cm y queremos dibujar uno semejante con un área de 600 cm². ¿Qué dimensiones tiene que tener?

\ex{3} Un triángulo rectángulo tiene catetos de 3 cm y 4 cm. Un triángulo semejante tiene 54 cm² de área. ¿Cuánto miden sus lados?


% =====================================================
\seccio{6}{Criterios de semejanza de triángulos}

\ex{1} ¿Son semejantes los triángulos que tienen estos lados? Si lo son, da la razón de semejanza.\begin{tasks}(2)
  \task 2, 3, 4 \ y \ 6, 9, 12
  \task 3, 4, 6 \ y \ 6, 8, 10
  \task 5, 5, 8 \ y \ 7,5, 7,5, 12
  \task 4, 6, 9 \ y \ 6, 9, 12
\end{tasks}\recorda{Dos triángulos son semejantes si tienen: \textbf{1)} dos ángulos iguales; \ o \textbf{2)} los tres lados proporcionales; \ o \textbf{3)} dos lados proporcionales y el ángulo comprendido igual. Para comparar lados, ordénalos de menor a mayor.}

\ex{1} ¿Son semejantes estos triángulos? Justifícalo.\begin{tasks}(1)
  \task Uno tiene ángulos de 35° y 75°, y el otro, de 75° y 70°.
  \task Uno tiene ángulos de 50° y 60°, y el otro, de 60° y 80°.
  \task Dos triángulos rectángulos, uno con un ángulo de 30° y el otro con un ángulo de 60°.
\end{tasks}

\ex{2} Los triángulos $ABC$ y $A'B'C'$ tienen $\hat A=\hat A'=50^\circ$, $\overline{AB}=4$, $\overline{AC}=6$, $\overline{A'B'}=6$ y $\overline{A'C'}=9$. ¿Son semejantes? Si $\overline{BC}=5$, ¿cuánto mide $\overline{B'C'}$?

\ex{2} Un triángulo tiene lados de 6, 8 y 10 cm. El lado pequeño de un triángulo semejante mide 9 cm. Calcula los otros dos lados y la razón entre las áreas.

\ex{2} ¿Son semejantes dos triángulos isósceles que tienen el ángulo desigual de 40°? ¿Y dos triángulos isósceles cualesquiera?

\ex{2} Los catetos de un triángulo rectángulo miden 6 cm y 8 cm. Calcula la hipotenusa y, usando triángulos semejantes, la altura sobre la hipotenusa.

\ex{3} Los catetos de un triángulo rectángulo $ABC$ (ángulo recto en $A$) miden $\overline{AB}=9$ cm y $\overline{AC}=12$ cm. La altura $AH$ divide la hipotenusa en dos segmentos. Calcula $\overline{BC}$, $\overline{BH}$, $\overline{HC}$ y $\overline{AH}$.

\ex{3} En un trapecio $ABCD$ las bases miden $\overline{AB}=12$ cm y $\overline{CD}=8$ cm, y las diagonales se cortan en $O$. Si $\overline{AO}=9$ cm, ¿cuánto mide $\overline{OC}$?


% =====================================================
\seccio{7}{Aplicaciones: alturas y distancias}

\ex{1} Un palo de 1 m proyecta una sombra de 1,25 m. En el mismo momento, un árbol proyecta una sombra de 10 m. ¿Qué altura tiene el árbol?\par\smallskip\begin{center}\begin{tikzpicture}[x=0.32cm,y=0.32cm,baseline=(current bounding box.center),line join=round]
\draw[muted] (-0.600,0.000) -- (17.000,0.000);
\draw[fill=accent!22,draw=accent,thick] (0.000,0.000) -- (0.000,2.000) -- (2.500,0.000) -- cycle;
\draw[thin] (0.688,0.000) -- (0.688,0.688) -- (0.000,0.688);
\draw[fill=teal!12,draw=teal,thick] (6.200,0.000) -- (6.200,8.000) -- (16.200,0.000) -- cycle;
\draw[thin] (6.888,0.000) -- (6.888,0.688) -- (6.200,0.688);
\draw[line width=2pt,draw=muted] (0.000,0.000) -- (0.000,2.000);
\draw[line width=2pt,draw=muted] (6.200,0.000) -- (6.200,8.000);
\draw[dashed,draw=muted] (0.000,2.000) -- (2.500,0.000);
\draw[dashed,draw=muted] (6.200,8.000) -- (16.200,0.000);
\node[font=\small,inner sep=2pt,left] at (0.000,1.000) {$1$\,m};
\node[font=\small,inner sep=2pt,below] at (1.250,0.000) {$1{,}25$\,m};
\node[font=\small,inner sep=2pt,left,text=accent] at (6.200,4.000) {$H$};
\node[font=\small,inner sep=2pt,below] at (11.200,0.000) {$10$\,m};
\end{tikzpicture}\end{center}\recorda{Busca dos \textbf{triángulos semejantes} (rayos de sol paralelos, un espejo, una visual que pasa por dos puntos…) y escribe la proporción entre lados homólogos: $\frac{\text{altura}}{\text{sombra}}=\frac{\text{altura}'}{\text{sombra}'}$.}

\ex{1} Un edificio de 24 m proyecta una sombra de 18 m. ¿Qué sombra proyecta, en el mismo momento, una persona de 1,6 m?

\ex{2} Marta tiene los ojos a 1,5 m del suelo. Pone un espejo $M$ en el suelo a 2 m de ella y ve reflejada la punta de una torre que está a 18 m del espejo. ¿Qué altura tiene la torre?\par\smallskip\begin{center}\begin{tikzpicture}[x=0.26cm,y=0.26cm,baseline=(current bounding box.center),line join=round]
\draw[muted] (-0.500,0.000) -- (21.000,0.000);
\draw[line width=2pt,draw=muted] (0.000,0.000) -- (0.000,1.500);
\draw[line width=2pt,draw=muted] (20.000,0.000) -- (20.000,13.500);
\draw[dashed,draw=muted] (0.000,1.500) -- (2.000,0.000);
\draw[dashed,draw=muted] (2.000,0.000) -- (20.000,13.500);
\fill (2.000,0.000) circle (1.5pt);
\node[font=\small,left] at (0.000,0.750) {$1{,}5$\,m};
\node[font=\small,below] at (1.000,-1.600) {$2$\,m};
\node[font=\small,below] at (11.000,-1.600) {$18$\,m};
\node[font=\small,right] at (20.000,6.750) {$H$};
\node[font=\small,below] at (2.000,0.000) {$M$};
\end{tikzpicture}\end{center}

\ex{2} Tales midió la gran pirámide con las sombras. Si en un momento un bastón de 2 m proyecta una sombra de 2 m y la sombra de la pirámide, contada desde el centro de la base, mide 146 m, ¿qué altura tiene la pirámide?

\ex{2} Una estaca de 1,5 m está clavada a 2 m de un punto $P$ del suelo. Desde $P$, la punta de la estaca y la punta de una torre que está a 30 m de $P$ quedan alineadas. ¿Qué altura tiene la torre?

\ex{3} Una persona de 1,6 m está a 3 m de una farola de 5 m de altura. ¿Cuánto mide su sombra?\par\smallskip\begin{center}\begin{tikzpicture}[x=0.55cm,y=0.55cm,baseline=(current bounding box.center),line join=round]
\draw[muted] (-0.500,0.000) -- (5.200,0.000);
\draw[line width=2pt,draw=muted] (0.000,0.000) -- (0.000,5.000);
\draw[line width=2pt,draw=muted] (3.000,0.000) -- (3.000,1.600);
\draw[dashed,draw=muted] (0.000,5.000) -- (4.412,0.000);
\fill (0.000,5.000) circle (1.5pt);
\node[font=\small,left] at (0.000,2.500) {$5$\,m};
\node[font=\small,left] at (3.000,0.800) {$1{,}6$\,m};
\node[font=\small,below] at (1.500,0.000) {$3$\,m};
\node[font=\small,below] at (3.706,0.000) {$s$};
\end{tikzpicture}\end{center}

\ex{3} Para medir la anchura $\overline{AB}$ de un río, desde $B$ (delante de un árbol $A$ de la otra orilla) caminamos 20 m perpendicularmente a $AB$ hasta $C$ y 8 m más hasta $D$. Desde $D$ nos alejamos del río, perpendicularmente, hasta el punto $E$ desde donde vemos $C$ y $A$ alineados. Si $\overline{DE}=6$ m, ¿qué anchura tiene el río?

\ex{3} Una persona con los ojos a 1,5 m del suelo está a 3 m de una tapia de 3 m. Ve la parte de arriba de la tapia alineada con la parte de arriba de un edificio que está 7 m detrás de la tapia. ¿Qué altura tiene el edificio?\par\smallskip\begin{center}\begin{tikzpicture}[x=0.45cm,y=0.45cm,baseline=(current bounding box.center),line join=round]
\draw[muted] (-0.500,0.000) -- (10.800,0.000);
\draw[line width=2pt,draw=muted] (0.000,0.000) -- (0.000,1.500);
\draw[line width=2pt,draw=muted] (3.000,0.000) -- (3.000,3.000);
\draw[line width=2pt,draw=muted] (10.000,0.000) -- (10.000,6.500);
\draw[dashed,draw=muted] (0.000,1.500) -- (10.000,6.500);
\draw[dashed,draw=muted] (0.000,1.500) -- (10.000,1.500);
\node[font=\small,left] at (0.000,0.750) {$1{,}5$};
\node[font=\small,right] at (3.000,1.500) {$3$\,m};
\node[font=\small,below] at (1.500,0.000) {$3$\,m};
\node[font=\small,below] at (6.500,0.000) {$7$\,m};
\node[font=\small,right] at (10.000,4.000) {$H$};
\end{tikzpicture}\end{center}


% =====================================================
\seccio{8}{Escalas. Áreas y volúmenes de figuras semejantes}

\ex{1} Calcula la distancia real:\begin{tasks}(3)
  \task 8 cm en un plano 1:50
  \task 6 cm en un mapa 1:25\,000
  \task 3,5 cm en un mapa 1:200\,000
\end{tasks}\recorda{Escala $1:n$: 1 cm del plano son $n$ cm reales. \ Real $=$ plano $\times\,n$; \ plano $=$ real $:\,n$. Si la razón de semejanza es $k$, las áreas se multiplican por $k^2$ y los volúmenes por $k^3$.}

\ex{1} ¿Qué longitud tendrá en el plano? \ a) Una pared de 12 m, en un plano 1:100. \ b) Una carretera de 45 km, en un mapa 1:1\,500\,000.

\ex{2} Calcula la escala: \ a) 5 cm del mapa representan 2 km; \ b) 3 cm del plano representan 150 m.

\ex{2} En el plano de un piso a escala 1:200, el comedor mide 3 cm × 2,5 cm. Calcula sus dimensiones reales y su superficie.

\ex{2} La escala gráfica de un mapa indica que 2 cm son 500 m. ¿Cuál es la escala numérica? ¿Qué distancia real hay entre dos puntos que están a 7 cm en el mapa?

\ex{2} Una maqueta de un coche está hecha a escala 1:18. Si el coche real mide 4,5 m, ¿cuánto mide la maqueta? El parabrisas de la maqueta tiene 40 cm². ¿Qué superficie tiene el real?

\ex{3} Dos botellas semejantes miden 15 cm y 30 cm de altura. La pequeña tiene una capacidad de 0,25 L y una etiqueta de 60 cm². ¿Qué capacidad tiene la grande? ¿Y qué superficie tiene su etiqueta?

\ex{3} En un mapa a escala 1:50\,000, un bosque ocupa 6 cm². ¿Qué superficie real tiene, en km² y en hectáreas?

\vfill
\begin{center}\color{muted}\small\faCheckCircle[regular]\ \ ¡Ánimo! Si te cuesta algún ejercicio, repasa los apuntes interactivos de la unidad.\end{center}

\end{document}
